\documentclass[12pt]{article}
\usepackage[utf8]{inputenc}
\usepackage[T1]{fontenc}
\usepackage{hyperref}
\title{Northern Tutchone Language Resources for K-12 Curriculum Development}
\author{Compiled by Claude Sonnet 5 (claude-sonnet-5) for Andr\'as Kornai}
\date{}
\begin{document}
\maketitle
\section{Introduction}
Northern Tutchone (ISO 639-3 ttm) is an endangered Athabaskan language of Yukon, Canada. A 2011 census count cited on Wikipedia gave 210 speakers. Several older print materials have been digitised, mainly through the Yukon Native Language Centre (YNLC). No current Northern Tutchone school curriculum was found. The school programme details in the results concerned Southern Tutchone.
\section{Primary Educational Resources}
\begin{itemize}
\item \emph{Northern Tutchone Language Lessons: Carmacks Dialect} (1998), a bilingual phrasebook with audio from speakers Eva Billy and Grace Wheeler. \url{https://wiki.ubc.ca/Documentation:RelLex/Northern_Tutchone_Language_Lessons:_Carmacks_Dialect}
\end{itemize}
\section{Linguistic and Reference Materials}
\begin{itemize}
\item \emph{The Selkirk Indian Language Noun Dictionary} (1977), English to Northern Tutchone. \url{https://wiki.ubc.ca/Documentation:RelLex/The_Selkirk_Indian_Language_Noun_Dictionary_(Northern_Tutchone_Athapascan)}
\item \emph{Mayo Indian Language: Noun Dictionary}, by Ritter, Hager and Peter, undated. \url{https://wiki.ubc.ca/Documentation:RelLex/Mayo_Indian_Language:_Noun_Dictionary}
\item NNDFN Dictionary, hosted by YNLC. \url{https://ynlc.ca/wp-content/uploads/NNDFN-Dictionary.pdf}
\item Omniglot, Tutchone. \url{https://www.omniglot.com/writing/tutchone.htm}
\item FreeLang Northern Tutchone wordlist. \url{https://www.freelang.net/dictionary/tutchone_northern.php}
\end{itemize}
\section{Religious and Community Resources}
None found beyond the YNLC materials.
\section{Digital and Audio Resources}
CBC reported in 2018 that over 10,000 pages of Yukon Indigenous language material, including Northern Tutchone, were published online and made searchable. \url{https://www.cbc.ca/1.4761056}
\section{Curriculum Development Context}
There are digitised word lists and a bilingual phrasebook, but no modern curriculum was found. A curriculum could build on the YNLC scans and on work with the remaining first-language speakers.
\end{document}